\section{Methodologie du projet}
La gestion de projet est une etape essentielle pour assurer la livraison
reussie d'une plateforme multi-locataire de receptionniste virtuelle.
Plusieurs methodologies de gestion de projet sont disponibles pour aider
les equipes a planifier, executer et adapter leur travail efficacement.
Parmi celles-ci, la methodologie Scrum est particulierement adaptee aux
projets qui combinent developpement d'IA conversationnelle, exigences
metier evolutives, et retours rapides des clients --- caracteristiques
qui definissent precisement le projet Nova~AI.

\subsection{La methode Scrum}
Scrum est aujourd'hui la methode Agile la plus populaire. L'approche Scrum
suit les principes du manifeste Agile, ce qui implique une participation
active du client tout au long du projet. Le terme \og scrum \fg{} provient
du vocabulaire sportif et designe la melee au rugby. Le principe de base
est que l'equipe avance comme une seule unite et reste prete a reorienter
le projet a mesure qu'il progresse, a l'image d'une melee ou les joueurs,
en groupe solidaire, s'adaptent collectivement au mouvement du ballon.

La methodologie Scrum repose sur des \emph{sprints}, qui sont des
intervalles de temps relativement courts, generalement de deux a quatre
semaines. A la fin de chaque sprint, l'equipe presente ce qu'elle a
ajoute au produit. La figure \ref{fig:scrum-cycle} resume le cycle de vie
d'un projet Scrum.

% ===========================================================================
% Cycle de vie d'un projet Scrum
% ===========================================================================
\begin{figure}[H]
\centering
\begin{tikzpicture}[
  font=\footnotesize,
  every node/.style={align=center},
  pb/.style    ={rectangle, rounded corners=2pt, draw=black!75,
                 fill=violet!15, minimum width=30mm, minimum height=14mm,
                 font=\scriptsize\bfseries},
  sb/.style    ={pb, fill=cyan!12, minimum width=30mm},
  ev/.style    ={pb, fill=orange!12, minimum width=24mm},
  artif/.style ={pb, fill=green!12, minimum width=24mm},
  arrow/.style ={->, thick, draw=black!75, >=Stealth},
  loop/.style  ={->, thick, draw=violet!70, >=Stealth, dashed},
]

% Linear flow left → right
\node[pb]    (pbacklog) at (0, 0) {Product\\Backlog};
\node[ev]    (planning) at (3.6, 0) {Sprint\\Planning};
\node[sb]    (sbacklog) at (7.0, 0) {Sprint\\Backlog};

% Sprint loop (in the middle)
\node[ev]    (daily) at (10.5, 1.6) {Daily\\Scrum};
\node[artif] (incr)  at (10.5, 0)   {Incrément};
\node[ev]    (review)at (10.5,-1.6) {Sprint\\Review};

% Final
\node[ev]    (retro) at (14.5, 0) {Sprint\\Retrospective};

% Linear arrows
\draw[arrow] (pbacklog) -- (planning);
\draw[arrow] (planning) -- (sbacklog);
\draw[arrow] (sbacklog) -- ++(2.0,0) |- (daily.west);

% Sprint inner cycle (24h loop daily-incrément-daily)
\draw[loop] (daily.south)  -- (incr.north);
\draw[loop] (incr.south)   -- (review.north);
\draw[loop] (review.west)  to[bend left=45]
   node[midway, fill=white, font=\scriptsize, inner sep=1pt]{2--4 sem.}
   (daily.west);

% From review/retrospective back to product backlog (refinement)
\draw[arrow] (incr.east)  -- (retro.west);
\draw[arrow] (retro.north) to[bend right=30]
   node[midway, fill=white, font=\scriptsize, inner sep=1pt]{nouveau sprint}
   (pbacklog.north);

% Caption notes
\node[font=\scriptsize\itshape, below=14mm of pbacklog, anchor=north, xshift=0mm] (legend1) {priorisé par le Product Owner};
\node[font=\scriptsize\itshape, below=14mm of sbacklog, anchor=north] (legend2) {engagement de l'équipe};
\node[font=\scriptsize\itshape, below=14mm of incr, anchor=north] (legend3) {potentiellement livrable};
\node[font=\scriptsize\itshape, below=14mm of retro, anchor=north] (legend4) {amélioration continue};

\end{tikzpicture}
\caption{Cycle de vie d'un projet Scrum : un Product Backlog priorisé alimente chaque Sprint Planning ; le Sprint Backlog devient l'engagement de l'équipe pour 2 à 4 semaines, rythmé par des Daily Scrums quotidiens, et clôturé par une Sprint Review (démo) et une Sprint Retrospective (amélioration continue).}
\label{fig:scrum-cycle}
\end{figure}


Pour le projet Nova~AI, nous avons structure notre processus Scrum afin
de prendre en compte trois flux de travail entrelaces : le developpement
back-end sur Supabase (schema de base de donnees, fonctions Edge,
souscriptions temps reel), le travail front-end sur le tableau de bord
multi-locataire et la PWA patient, et l'ingenierie de prompts pour le
receptionniste IA qui gere les conversations WhatsApp en arabe, francais
et anglais.

\subsubsection{Roles Scrum}
\begin{itemize}
  \item \textbf{L'equipe.} Responsable de transformer les exigences
        definies par le Product Owner en fonctionnalites utilisables.
        Dans notre projet, l'equipe est full-stack et combine ingenierie
        front-end sur le tableau de bord, ingenierie back-end sur Postgres
        et les fonctions Edge Deno, ingenierie IA pour le classifieur
        d'intention base sur Claude et le verificateur de securite des
        prescriptions, ainsi que conception UI/UX pour la page d'accueil
        au theme \emph{anti-gravite} et l'interface de la PWA.
  \item \textbf{Le Scrum Master.} Responsable de la comprehension, du
        respect et de l'organisation de la methodologie Scrum. Il garantit
        que les iterations de prompts IA, les migrations de schema et les
        evolutions du tableau de bord progressent harmonieusement et
        conformement aux principes de Scrum, en levant les obstacles tels
        que les delais de verification de l'API Meta WhatsApp ou les
        contraintes de quota d'ElevenLabs.
  \item \textbf{Le Product Owner.} Represente generalement le client et
        porte la vision du produit a creer. Pour notre projet, le role de
        Product Owner agrege les perspectives de trois personae cibles :
        un proprietaire de clinique tunisien, un responsable de
        laboratoire d'analyses, et un medecin independant. Leurs points
        de douleur recurrents (rendez-vous manques, prescriptions
        manuscrites, prises de rendez-vous hors heures) ont oriente la
        priorisation du Product Backlog.
\end{itemize}

\subsubsection{Evenements Scrum}
\begin{itemize}
  \item \textbf{Sprint Planning.} Les taches a accomplir durant le sprint
        sont determinees par l'ensemble de l'equipe Scrum lors de la
        reunion de planification. Pour Nova~AI, cette planification couvre
        a la fois le travail produit (un nouvel onglet du tableau de bord,
        une migration Postgres, un gestionnaire de reponse a un bouton
        WhatsApp) et le travail IA (raffinement de la machine d'etats de
        prise de rendez-vous, extension du prompt de securite des
        medicaments).
  \item \textbf{Sprint Review.} Il s'agit de l'evaluation du sprint
        accompli. Une fois le sprint termine, l'equipe et les parties
        prenantes se reunissent pour valider ce qui a ete realise. Pour
        notre projet, cela inclut la demonstration des nouveaux onglets du
        tableau de bord, l'execution d'un parcours de prise de rendez-vous
        de bout en bout sur un vrai numero Meta, et la revue des
        transcriptions de l'IA sur les cas limites.
  \item \textbf{Sprint Retrospective.} Vise a s'adapter aux changements et
        a ameliorer continuellement le processus de mise en oeuvre.
        L'equipe utilise la retrospective pour examiner le sprint termine
        et determiner ce qui a bien fonctionne et ce qui doit etre
        ameliore. Cette pratique a ete particulierement precieuse pour
        affiner la strategie de prompt IA apres avoir observe des
        mauvaises classifications en conditions reelles, et pour
        renforcer les politiques de Row-Level Security multi-locataire
        apres chaque nouvelle fonctionnalite.
  \item \textbf{Daily Scrum.} Cette reunion quotidienne de quinze minutes
        est tres importante. Son objet est de suivre la progression
        quotidienne du sprint. Pour notre projet, le stand-up quotidien
        couvre l'avancement des composants du tableau de bord, des
        deploiements de fonctions Edge, des revisions de prompts IA et
        des tests de migrations Supabase.
\end{itemize}

\subsubsection{Artefacts Scrum}
\begin{itemize}
  \item \textbf{Le Product Backlog.} Une liste priorisee des exigences
        initiales du client pour le produit a creer. Le Product Owner est
        responsable du Product Backlog. Notre Product Backlog couvre les
        user stories pour la prise de rendez-vous, les rappels, le
        pipeline de laboratoire, le remplissage de creneaux a la demande,
        l'apprentissage des durees de consultation, la securite des
        medicaments, les profils de vehicules, les surprises
        d'anniversaire, le check-in QR, la file d'attente en temps reel,
        et l'application web progressive pour les patients.
  \item \textbf{Le Sprint Backlog.} Le plan detaille pour atteindre
        l'objectif du sprint, defini lors de la reunion de planification.
        Pour notre projet, chaque Sprint Backlog contenait un melange de
        migrations de base de donnees, de taches de fonctions Edge, de
        travaux d'IHM et d'experiences d'ingenierie de prompts ---
        typiquement de trois a six user stories dimensionnees pour tenir
        dans un sprint de deux semaines.
\end{itemize}

\subsection{Justification du choix de la methodologie Agile Scrum}
La methodologie Agile Scrum est particulierement adaptee au projet
Nova~AI en raison de sa flexibilite intrinseque et de sa nature
iterative, qui s'alignent bien avec les exigences evolutives d'une
plateforme IA multi-locataire desservant des entreprises heterogenes
(cliniques, laboratoires, salons, garages, medecins independants).
Comme le systeme inclut des fonctionnalites avancees telles qu'un
receptionniste IA multi-tour, une messagerie WhatsApp en temps reel,
une OCR de prescriptions par modele de vision, des durees de
consultation auto-apprises, du remplissage de creneaux a la demande,
et une file d'attente en temps reel, Scrum a permis a l'equipe de
livrer des increments fonctionnels de maniere fiable tout en
accommodant les priorites changeantes.

L'accent mis par la methodologie sur la collaboration trans-fonctionnelle
a favorise une integration sans heurts entre la logique front-end
(l'IHM du tableau de bord au theme anti-gravite), les services back-end
(Postgres Supabase, fonctions Edge, Storage, Realtime) et les composants
IA (Claude Haiku pour la classification d'intention, Claude Sonnet pour
l'OCR des images de prescriptions, ElevenLabs en option pour le clonage
vocal). De plus, son cadre ouvert avec des stand-ups quotidiens et des
sprint reviews garantit que les parties prenantes restent engagees tout
au long du processus de developpement.

La methodologie a ete particulierement precieuse en cela qu'au fur et a
mesure que le projet s'est etendu pour englober de nouveaux types
d'entreprises et de nouvelles capacites --- profils de vehicules pour
garages, gestion de file d'attente pour laboratoires, verificateur
d'interactions medicamenteuses --- elle a soutenu une amelioration
continue basee sur les retours utilisateurs reels et sur l'evolution
des besoins du marche dans le contexte tunisien.

Cette methode a ete tres benefique pour notre projet, car elle nous a
permis de mettre en oeuvre efficacement la prise de rendez-vous assistee
par IA, la gestion temps reel des rendez-vous, et les flux metiers
intelligents pour les petites et moyennes entreprises tunisiennes, tout
en restant centre sur les objectifs metier et en adressant les defis
techniques d'un developpement SaaS multi-locataire.
